Ein zweiatomiges Gas hat bei Zimmertemperatur $f = \fzO$ Freiheitsgrade (drei der Translation, zwei der Rotation). Mit $C_v = \frac{f}{2}R$ und $C_p = C_v + R = \frac{f+2}{2}R$ (FoTaBe S.~171) ist
\begin{align}
 \kappa &= \frac{C_p}{C_v} = \kapF = \frac{\fzO+2}{\fzO} = \kap \approx \result{\kapP{0}{2}}
\end{align}
Die Tabelle (FoTaBe S.~191) bestätigt das, z.\,B. Stickstoff $\kappa = \kNiO$, Sauerstoff $\kappa = \kOxO$.
